% =====================================================================
% CAP. 11 - NIVEL 4
% =====================================================================
\blocodesafio

\begin{questao}[4]{}
Em laboratório, um circuito RC com $R=\SI{5}{\kilo\ohm}$ e
$C=\SI{10}{\micro\farad}$ deveria apresentar $\tau=\SI{50}{\milli\second}$.
Entretanto, o ajuste exponencial do osciloscópio fornece
$\tau_{\mathrm{med}}=\SI{80}{\milli\second}$. A capacitância medida está
próxima do valor nominal. Estime a resistência efetiva vista pelo capacitor e proponha
uma causa elétrica plausível para a discrepância.
\resposta{$R_{\mathrm{eq}}\approx\SI{8}{\kilo\ohm}$; há cerca de \SI{3}{\kilo\ohm}
adicionais na trajetória equivalente}
\resolucao{$R_{\mathrm{eq}}=\tau/C=0{,}08/(\pot{10}{-6})=\SI{8}{\kilo\ohm}$.
Como o resistor nominal é de \SI{5}{\kilo\ohm}, há uma resistência
adicional de aproximadamente \SI{3}{\kilo\ohm}. Essa diferença pode decorrer da resistência de saída da fonte, de uma
rede de polarização não considerada, de um resistor em série, de uma chave ou
sensor, ou de uma topologia
que não foi corretamente reduzida ao equivalente de Thévenin.}
\end{questao}

\begin{questao}[4]{}
Um capacitor de $\SI{470}{\micro\farad}$ será carregado por uma fonte de
$\SI{24}{\volt}$. Deseja-se limitar a corrente inicial a
$\SI{100}{\milli\ampere}$ e atingir pelo menos \SI{95}{\percent} da tensão final
em menos de $\SI{0.5}{\second}$. Projete o resistor mínimo em série, verifique o
tempo de carga e estime a energia total dissipada nele durante uma carga completa.
\resposta{$R=\SI{240}{\ohm}$; $t_{95}\approx\SI{0.338}{\second}$; $E_R\approx\SI{0.135}{\joule}$}
\resolucao{Pela corrente inicial, $R\ge24/0{,}1=\SI{240}{\ohm}$. Com esse valor,
$\tau=RC=240\cdot\pot{470}{-6}=\SI{112.8}{\milli\second}$. Para \SI{95}{\percent},
$t=-\tau\ln0{,}05\approx\SI{0.338}{\second}$, atendendo à especificação. A
energia dissipada no resistor durante a carga ideal por fonte de tensão é igual à
energia final do capacitor:
$E_R=\tfrac12CV^2\approx\SI{0.135}{\joule}$. A potência inicial é
$V^2/R=\SI{2.4}{\watt}$, importante para a escolha da tecnologia do resistor.}
\end{questao}

\begin{questao}[4]{}
Uma bobina de relé tem $L=\SI{120}{\milli\henry}$ e resistência de cobre
$\SI{60}{\ohm}$, alimentada em $\SI{24}{\volt}$ até o regime permanente.
Compare duas estratégias de desligamento: (A) diodo de roda livre ideal; (B)
resistor de clamp de $\SI{180}{\ohm}$ no caminho de descarga. Determine a energia
inicial, as constantes de tempo aproximadas e a tensão inicial no resistor de
clamp. Explique o compromisso de projeto.
\resposta{$I_0=\SI{0.4}{\ampere}$; $E_L=\SI{9.6}{\milli\joule}$; $\tau_A\approx\SI{2}{\milli\second}$;
$\tau_B\approx\SI{0.5}{\milli\second}$; $v_{R,\mathrm{clamp}}(0^+)\approx\SI{72}{\volt}$}
\resolucao{Em regime, $I_0=24/60=\SI{0.4}{\ampere}$ e
$E_L=\tfrac12LI_0^2=\SI{9.6}{\milli\joule}$. Com o diodo ideal, a resistência no caminho de descarga é aproximadamente
a resistência do enrolamento da bobina: $\tau_A=L/60=\SI{2}{\milli\second}$. Com o
resistor adicional, $R_{\mathrm{tot}}=60+180=\SI{240}{\ohm}$ e
$\tau_B=L/240=\SI{0.5}{\milli\second}$. A tensão inicial no resistor externo é
$0{,}4\cdot180=\SI{72}{\volt}$. O clamp resistivo desliga o relé mais rápido,
mas produz tensão maior; o diodo reduz a sobretensão, porém prolonga a corrente.}
\end{questao}

\begin{questao}[4]{}
Uma oscilação amortecida medida em um circuito RLC apresenta período entre picos
de aproximadamente $\SI{4}{\milli\second}$ e a envoltória cai à metade a cada
$\SI{10}{\milli\second}$. O capacitor é $C=\SI{10}{\micro\farad}$. Estime
$\alpha$, $\omega_d$, $\omega_0$, $\zeta$, $L$ e a resistência equivalente em série.

\begin{figure}[H]
\centering
\begin{tikzpicture}
\begin{axis}[
 width=0.88\linewidth,height=5.2cm,
 xlabel={$t$ (ms)},ylabel={$v(t)$},xmin=0,xmax=30,
 axis lines=middle,grid=both,samples=300,
 tick label style={font=\scriptsize},label style={font=\small\sffamily}]
\addplot[thick,azulprof,domain=0:30] {10*exp(-0.0693147*x)*cos(deg(1.5707963*x))};
\addplot[dashed,laranja,domain=0:30] {10*exp(-0.0693147*x)};
\addplot[dashed,laranja,domain=0:30] {-10*exp(-0.0693147*x)};
\end{axis}
\end{tikzpicture}
\caption{Exemplo de oscilação amortecida e sua envoltória exponencial.}
\end{figure}

\resposta{$\alpha\approx\SI{69.3}{\per\second}$; $\omega_d\approx\SI{1571}{\radian\per\second}$;
$\omega_0\approx\SI{1572}{\radian\per\second}$; $\zeta\approx0{,}044$;
$L\approx\SI{40.45}{\milli\henry}$; $R\approx\SI{5.61}{\ohm}$}
\resolucao{Da envoltória, $e^{-\alpha(0{,}01)}=1/2$, então
$\alpha=\ln2/0{,}01\approx\SI{69.3}{\per\second}$. Do período,
$\omega_d=2\pi/0{,}004\approx\SI{1571}{\radian\per\second}$. Logo
$\omega_0=\sqrt{\omega_d^2+\alpha^2}\approx\SI{1572}{\radian\per\second}$ e
$\zeta\approx0{,}044$. Como $\omega_0^2=1/(LC)$,
$L\approx1/(C\omega_0^2)\approx\SI{40.45}{\milli\henry}$. Por fim,
$R=2\alpha L\approx\SI{5.61}{\ohm}$.}
\end{questao}

\begin{questao}[4]{}
Um estágio RLC em série usa $L=\SI{20}{\milli\henry}$ e
$C=\SI{10}{\micro\farad}$. O requisito é obter a resposta mais rápida possível
sem oscilação. Determine o resistor ideal e compare qualitativamente o uso de
$\SI{45}{\ohm}$ e $\SI{180}{\ohm}$.
\resposta{$R_{\mathrm{crit}}\approx\SI{89.4}{\ohm}$; \SI{45}{\ohm} é subamortecido;
\SI{180}{\ohm} é superamortecido}
\resolucao{$R_{\mathrm{crit}}=2\sqrt{L/C}=2\sqrt{0{,}02/\pot{10}{-6}}
\approx\SI{89.4}{\ohm}$. Com \SI{45}{\ohm}, $R<R_{\mathrm{crit}}$ e haverá
oscilações amortecidas (ringing) e sobressinal. Com \SI{180}{\ohm}, $R>R_{\mathrm{crit}}$ e a resposta será
monotônica, porém mais lenta. O amortecimento crítico é o limite de resposta
monotônica mais rápida para o modelo ideal.}
\end{questao}

\begin{questao}[4]{}
A tabela mostra a resposta de um RC inicialmente descarregado cuja tensão final é
$\SI{10}{\volt}$.

\begin{table}[H]
\centering\small\sffamily
\begin{tabular}{@{}c c@{}}
\toprule
$t$ & $v_C(t)$\\
\midrule
\SI{20}{\milli\second} & \SI{6.32}{\volt}\\
\SI{40}{\milli\second} & \SI{8.65}{\volt}\\
\bottomrule
\end{tabular}
\end{table}

Verifique se os dados são compatíveis com uma única constante de tempo e estime
$R$ se $C=\SI{4.7}{\micro\farad}$.
\resposta{$\tau\approx\SI{20}{\milli\second}$; $R\approx\SI{4.26}{\kilo\ohm}$}
\resolucao{Em $t=\tau$, uma carga RC atinge \SI{63.2}{\percent}, exatamente o
primeiro ponto. Em $2\tau$, atinge $1-e^{-2}=0{,}8647$, coerente com
\SI{8.65}{\volt}. Logo $\tau\approx\SI{20}{\milli\second}$. Assim
$R=\tau/C=0{,}02/(\pot{4.7}{-6})\approx\SI{4.26}{\kilo\ohm}$.}
\end{questao}

\begin{questao}[4]{}
Após uma comutação, um circuito RLC em série fica sem fonte independente e possui
$R=\SI{20}{\ohm}$, $L=\SI{50}{\milli\henry}$ e
$C=\SI{100}{\micro\farad}$. Adote $v_C(0^+)=\SI{5}{\volt}$ e
$i(0^+)=\SI{0.1}{\ampere}$ entrando no terminal positivo do capacitor. Determine
a expressão de $v_C(t)$.
\resposta{$v_C(t)=e^{-200t}[5\cos(400t)+5\sin(400t)]\ \si{\volt}$}
\resolucao{$\alpha=R/(2L)=200$, $\omega_0=1/\sqrt{LC}\approx447{,}2$ e
$\omega_d=\sqrt{\omega_0^2-\alpha^2}=400$. Portanto
$v_C=e^{-200t}(A\cos400t+B\sin400t)$. De $v_C(0)=5$, $A=5$. Como
$i=C\dot v_C$, $\dot v_C(0)=0{,}1/(\pot{100}{-6})=\SI{1000}{\volt\per\second}$.
Então $\dot v_C(0)=-200A+400B=1000$, resultando em $B=5$.}
\end{questao}

\begin{questao}[4]{}
Projete uma rede de temporização para uma entrada de $\SI{3.3}{\volt}$ e
$C=\SI{1}{\micro\farad}$. Na carga, a tensão deve alcançar
$\SI{2}{\volt}$ em aproximadamente $\SI{8}{\milli\second}$. No desligamento,
deve cair de $\SI{3.3}{\volt}$ para menos de $\SI{0.5}{\volt}$ em
$\SI{5}{\milli\second}$. Mostre que um único resistor não atende bem às duas
metas e proponha resistores independentes de carga e descarga usando um diodo.
\resposta{$R_{\mathrm{carga}}\approx\SI{8.59}{\kilo\ohm}$; $R_{\mathrm{descarga}}\lesssim\SI{2.65}{\kilo\ohm}$;
valores práticos: \SI{8.2}{\kilo\ohm} e \SI{2.4}{\kilo\ohm}}
\resolucao{Na carga,
$2=3{,}3(1-e^{-8\mathrm{ms}/RC})$, fornecendo
$R\approx\SI{8.59}{\kilo\ohm}$. Se esse mesmo resistor fosse usado na descarga,
o tempo até \SI{0.5}{\volt} seria $RC\ln(3{,}3/0{,}5)\approx\SI{16.2}{\milli\second}$,
acima do limite. Para a descarga em no máximo \SI{5}{\milli\second},
$R\le t/[C\ln(3{,}3/0{,}5)]\approx\SI{2.65}{\kilo\ohm}$. Um diodo permite
separar os caminhos; \SI{8.2}{\kilo\ohm} e \SI{2.4}{\kilo\ohm} são escolhas
comerciais coerentes.}
\end{questao}

\begin{questao}[4]{}
Para $L=\SI{20}{\milli\henry}$ e $C=\SI{10}{\micro\farad}$, compare os três
resistores da tabela e escolha o mais adequado quando se deseja resposta rápida
sem sobressinal.

\begin{table}[H]
\centering\small\sffamily
\begin{tabular}{@{}c c c@{}}
\toprule
\textbf{Opção} & $R$ & \textbf{Observação a determinar}\\
\midrule
A & \SI{40}{\ohm} & ?\\
B & \SI{89.4}{\ohm} & ?\\
C & \SI{180}{\ohm} & ?\\
\bottomrule
\end{tabular}
\end{table}

\resposta{A: $\zeta\approx0{,}447$, subamortecido; B: crítico; C: $\zeta\approx2{,}01$, superamortecido; escolher B}
\resolucao{Como $R_{\mathrm{crit}}\approx\SI{89.4}{\ohm}$, no RLC em série vale
$\zeta=R/R_{\mathrm{crit}}$. Assim A fornece $\zeta\approx0{,}447$ e oscila; B
fica praticamente em $\zeta=1$; C fornece $\zeta\approx2{,}01$ e é lento por
superamortecimento. Para rapidez sem sobressinal, B é a escolha adequada.}
\end{questao}

\begin{questao}[4]{}
Um projeto prevê $R=\SI{20}{\ohm}$, $L=\SI{0.1}{\henry}$ e
$C=\SI{100}{\micro\farad}$, portanto deveria ser subamortecido. Contudo, o ajuste
da resposta medida indica dois polos reais em aproximadamente
$-\SI{112.7}{\per\second}$ e $-\SI{887.3}{\per\second}$. Admitindo $L$ e $C$
corretos, estime a resistência efetiva em série e diagnostique a discrepância.
\resposta{$R_{\mathrm{efetivo}}\approx\SI{100}{\ohm}$; há cerca de \SI{80}{\ohm}
de resistência em série não contabilizada}
\resolucao{Para $s^2+(R/L)s+1/(LC)=0$, a soma dos polos é $-R/L$. A soma medida é
$-1000\ \si{\per\second}$; com $L=\SI{0.1}{\henry}$,
$R=1000L=\SI{100}{\ohm}$. O produto dos polos é aproximadamente
$100000\ \si{\per\second\squared}$, coerente com $1/(LC)$, confirmando que $L$ e
$C$ são plausíveis. Portanto existe cerca de \SI{80}{\ohm} extra, por ESR,
resistência do indutor, resistor de proteção, fonte ou fiação.}
\end{questao}

\begin{questao}[4]{}
Um barramento CC de $\SI{400}{\volt}$ possui capacitor de
$\SI{470}{\micro\farad}$. Por segurança, após desligar a fonte deseja-se reduzir
a tensão para no máximo $\SI{50}{\volt}$ em $\SI{60}{\second}$. Dimensione o
maior resistor de descarga admissível, escolha um valor comercial abaixo dele,
estime a potência inicial e a energia que será dissipada.
\resposta{$R_{\max}\approx\SI{61.4}{\kilo\ohm}$; escolha \SI{56}{\kilo\ohm};
$t\approx\SI{54.7}{\second}$; $P_0\approx\SI{2.86}{\watt}$; $E\approx\SI{37.6}{\joule}$}
\resolucao{Na descarga, $50=400e^{-t/RC}$. Assim
$R=t/[C\ln(400/50)]\approx\SI{61.4}{\kilo\ohm}$. Um valor de
\SI{56}{\kilo\ohm} fornece $t=RC\ln8\approx\SI{54.7}{\second}$. A potência
inicial é $P_0=V^2/R=400^2/56000\approx\SI{2.86}{\watt}$, portanto é prudente
usar resistor com margem de potência e tensão. A energia inicialmente armazenada,
e posteriormente dissipada, é $E=\tfrac12CV^2\approx\SI{37.6}{\joule}$.}
\end{questao}

\gabarito[Nível 4]
